\chapter{Módulo de Alineación Masiva}
\label{ch:batch}
% ==============================================================

\section{Motivación}

El flujo anterior validaba fotogramas de forma \textbf{unitaria}: una llamada API por imagen,
sin visión global del lote. Esto presentaba tres problemas:

\begin{enumerate}[leftmargin=1.5em]
  \item \textbf{Escalabilidad}: con 24 imágenes/día por 6 cámaras, revisar manualmente cada
        alineación era inviable.
  \item \textbf{Detección de deriva}: sin comparativa global era difícil detectar si una
        cámara había girado ligeramente con el tiempo.
  \item \textbf{Objetos dinámicos}: SIFT sin restricciones detectaba gaviotas y olas como
        features, produciendo homografías sesgadas.
\end{enumerate}

\section{Arquitectura de estados}

El módulo implementa una \textbf{máquina de estados no-destructiva} de 5 fases:

\begin{figure}[H]
\centering
\begin{tikzpicture}[node distance=0.6cm and 1.8cm]
  \node[block, fill=gray!8] (idle) {idle};
  \node[block, fill=SkyBlue!15, right=of idle] (config) {configuring};
  \node[block, fill=Orange!15, right=of config] (proc) {processing};
  \node[block, fill=OliveGreen!15, right=of proc] (review) {reviewing};
  \node[block, fill=NavyBlue!15, right=of review] (commit) {committed};

  \draw[arrow] (idle) -- (config) node[midway,above,label]{start};
  \draw[arrow] (config) -- (proc) node[midway,above,label]{launch};
  \draw[arrow] (proc) -- (review) node[midway,above,label]{ready};
  \draw[arrow] (review) -- (commit) node[midway,above,label]{commit};

  % discard arco
  \draw[arrow, bend right=30, color=BrickRed!60]
    (review.south) to node[below,label]{\textcolor{BrickRed}{discard}} (idle.south);
  \draw[arrow, bend right=25, color=BrickRed!60]
    (proc.south) to node[below,label,pos=0.3]{\textcolor{BrickRed}{discard}} (idle.south);
\end{tikzpicture}
\caption{Máquina de estados del módulo BatchAlignment}
\label{fig:state-machine}
\end{figure}

\begin{infobox}[Principio de no-destructividad]
Todo el trabajo intermedio vive en \texttt{/tmp/cv\_lit\_batch/\{job\_id\}/}.
Los originales de cada cámara \textbf{nunca} se modifican.
La escritura permanente ocurre \textbf{únicamente} al llamar al commit explícito.
\end{infobox}

\subsection{Estructura del directorio temporal por job}

\begin{codebox}[]
\begin{minted}{text}
/tmp/cv_lit_batch/{job_id}/
  aligned/    ← imágenes con warpPerspective aplicado (staging)
  diff/       ← mapas de delta por imagen
  blend/      ← blend 50/50 en escala de grises
\end{minted}
\end{codebox}

\section{Pipeline de alineación: algoritmo}

\subsection{Parámetros configurables}

\begin{table}[H]
\centering
\caption{Parámetros del pipeline de alineación (\texttt{batch\_alignment.py})}
\label{tab:align-params}
\begin{tabular}{L{4cm} C{2.5cm} L{6cm}}
\toprule
\textbf{Constante} & \textbf{Valor} & \textbf{Descripción} \\
\midrule
\texttt{SIFT\_FEATURES}   & 3\,000 & Keypoints máximos por imagen \\
\texttt{LOWE\_RATIO}      & 0.75   & Ratio test de Lowe para filtrar matches \\
\texttt{MIN\_INLIERS}     & 15     & Inliers mínimos RANSAC (con máscara) \\
\texttt{MIN\_INLIERS\_FULL} & 30   & Inliers mínimos sin máscara (fallback) \\
\texttt{RANSAC\_THRESH}   & 5.0~px & Umbral de reproyección RANSAC \\
\bottomrule
\end{tabular}
\end{table}

\subsection{Estrategia de dos pasadas}

Para combinar la robustez de las máscaras con la resiliencia ante zonas escasas de features:

\begin{codebox}[Lógica de dos pasadas (pseudocódigo)]
\begin{minted}{python}
# Pasada 1: con máscara de zonas estables
H, inliers, n_good, fail = _try_align(gray, ref_kp, ref_des,
                                       mask=feat_mask,
                                       threshold=MIN_INLIERS)  # 15
if inliers >= MIN_INLIERS:
    result.used_mask = True
    return H, inliers

# Pasada 2 (fallback): toda la imagen, umbral más alto
H, inliers, n_good, fail = _try_align(gray, ref_kp_full, ref_des_full,
                                       mask=None,
                                       threshold=MIN_INLIERS_FULL)  # 30
result.used_mask = False
\end{minted}
\end{codebox}

\subsection{Flujo completo del pipeline}

\begin{enumerate}[leftmargin=1.5em]
  \item Carga imagen base → extrae keypoints con SIFT (con y sin máscara).
  \item Por cada imagen del lote (tarea de fondo asíncrona):
    \begin{enumerate}[leftmargin=1em]
      \item Detecta features → matching FLANN con ratio test de Lowe (0.75).
      \item RANSAC → homografía $\mathbf{H}$ (dos pasadas).
      \item \texttt{cv2.warpPerspective} → imagen alineada → \texttt{aligned/}.
      \item \texttt{\_compute\_diff\_map} → mapa de delta → \texttt{diff/}.
      \item \texttt{\_blend\_grayscale} → blend 50/50 → \texttt{blend/}.
      \item \texttt{\_compute\_mean\_shift} → métrica escalar $\delta_{\text{px}}$.
      \item Actualiza \texttt{job.progress\_current}.
    \end{enumerate}
  \item \verb|job.status = "ready"| al finalizar.
\end{enumerate}

\section{Renderizado del mapa de delta}

El mapa de delta es una imagen compuesta de tres capas superpuestas que proporciona tanto
información visual intuitiva como métricas cuantitativas.

\subsection{Capa 1 — Heatmap de diferencia absoluta}

\begin{equation}
  \Delta_\text{abs}(i,j) = \frac{1}{3}\sum_{c \in \{R,G,B\}} |I_\text{ref}(i,j,c) - I_\text{ali}(i,j,c)|
\end{equation}

Se normaliza y se aplica \texttt{COLORMAP\_HOT} de OpenCV:

\begin{table}[H]
\centering
\caption{Leyenda del heatmap HOT}
\begin{tabular}{C{2cm} L{8cm}}
\toprule
\textbf{Color} & \textbf{Significado} \\
\midrule
Negro & $\Delta = 0$ — alineación perfecta \\
Rojo/naranja & Error de intensidad moderado \\
Amarillo & Error alto \\
Blanco & Error severo — zona mal alineada \\
\bottomrule
\end{tabular}
\end{table}

\subsection{Capa 2 — Divergencia de bordes (overlay cian)}

\begin{codebox}[]
\begin{minted}{python}
edges_ref = cv2.Canny(ref_gray, 40, 120)
edges_ali = cv2.Canny(aligned_gray, 40, 120)
edge_diff = cv2.dilate(cv2.absdiff(edges_ref, edges_ali), kernel_5x5)
composite[edge_diff > 0] = (255, 255, 0)  # cian en BGR
\end{minted}
\end{codebox}

Los bordes estructurales que \textbf{no coinciden} entre referencia y alineada aparecen en
\textbf{cian}, identificando exactamente qué objeto se desplazó y en qué dirección.

\subsection{Capa 3 — Métrica escalar $\delta_{\text{px}}$ (flujo óptico Farneback)}

\begin{equation}
  \delta_{\text{px}} = \text{mediana}\left(
    \sqrt{v_x(i,j)^2 + v_y(i,j)^2}
  \right)
\end{equation}

donde $(v_x, v_y)$ es el campo de flujo óptico residual post-warp calculado con Farneback.
Mide cuántos píxeles de desplazamiento sistemático quedan tras la homografía.

\textbf{Umbral de alerta} en la UI: $\delta_{\text{px}} > 3$.

\section{API REST del módulo}

\begin{table}[H]
\centering
\caption{Endpoints del router \texttt{/api/batch/}}
\label{tab:api-batch}
\begin{tabular}{C{1.5cm} L{5cm} L{6cm}}
\toprule
\textbf{Método} & \textbf{Ruta} & \textbf{Descripción} \\
\midrule
\texttt{POST}  & \texttt{/start}                    & Encola job + lanza pipeline en background \\
\texttt{GET}   & \texttt{/\{id\}/status}             & Polling de progreso (fase processing) \\
\texttt{GET}   & \texttt{/\{id\}/results}            & Lista completa de métricas (fase ready) \\
\texttt{GET}   & \texttt{/\{id\}/preview/original/\{fn\}} & Imagen original sin transformar \\
\texttt{GET}   & \texttt{/\{id\}/preview/aligned/\{fn\}} & Imagen alineada en staging \\
\texttt{GET}   & \texttt{/\{id\}/preview/diff/\{fn\}} & Mapa de delta compuesto \\
\texttt{GET}   & \texttt{/\{id\}/preview/blend/\{fn\}} & Blend 50/50 en escala de grises \\
\texttt{PATCH} & \texttt{/\{id\}/approve/\{fn\}}    & Override de aprobación individual \\
\texttt{POST}  & \texttt{/\{id\}/commit}             & Two-phase commit → \texttt{CAM\_X/aligned/} \\
\texttt{DELETE}& \texttt{/\{id\}}                   & Descarta job y elimina temp dir \\
\bottomrule
\end{tabular}
\end{table}

\section{Modelos de dominio}

\begin{codebox}[Dataclasses principales (batch\_alignment.py)]
\begin{minted}{python}
@dataclass
class AlignmentResult:
    filename:      str
    status:        str       # "ok" | "failed" | "reference"
    inliers:       int       # inliers RANSAC
    mean_shift_px: float     # mediana flujo óptico residual (px)
    H:             Optional[List]  # homografía 3×3 como lista plana
    approved:      bool = True     # override de aprobación por imagen
    fail_reason:   Optional[str] = None  # no_features|low_matches|ransac_failed
    used_mask:     bool = True     # True = pasada con máscara activa

@dataclass
class BatchJob:
    job_id:           str
    cam_id:           int
    base_filename:    str
    image_filenames:  List[str]
    status:           str    # pending|processing|ready|committed|failed|discarded
    results:          Dict[str, AlignmentResult]
    progress_current: int
    progress_total:   int
\end{minted}
\end{codebox}

\section{Flujo de commit (two-phase write)}

\begin{figure}[H]
\centering
\begin{tikzpicture}[node distance=0.7cm and 2cm]
  \node[block, fill=OliveGreen!10] (review) {reviewing};
  \node[block, fill=NavyBlue!10, right=3.5cm of review] (staging)
    {\texttt{/tmp/.../aligned/}\\{\tiny staging no-destructivo}};
  \node[block, fill=BrickRed!8, below=of staging] (disk)
    {\texttt{data/CAM\_X/aligned/}\\{\tiny escritura permanente}};
  \node[block, fill=gray!8, below=of review] (discard)
    {\texttt{shutil.rmtree(temp\_dir)}\\{\tiny discard — cero side-effects}};

  \draw[arrow] (review) -- (staging) node[midway,above,label]{commit (approved=True)};
  \draw[arrow] (staging) -- (disk) node[midway,right,label]{\texttt{shutil.copy2}};
  \draw[arrow, color=BrickRed!60] (review) -- (discard) node[midway,left,label]{discard};
\end{tikzpicture}
\caption{Flujo de commit no-destructivo}
\end{figure}

% ==============================================================
